\documentclass[11pt]{article}

% ============================================================
%  Research Technical Note
%  Topic: time_block_strategy
%  AR(2) temporal pooling in a spatial skew-t model:
%  diagnosis of a null result and a time-block holdout redesign.
% ============================================================

% ---- Packages ----
\usepackage{amsmath, amssymb, amsthm}
\usepackage{mathtools}
\usepackage{bm}
\usepackage[margin=1in]{geometry}
\usepackage{hyperref}
\usepackage{booktabs}
\usepackage{array}
\usepackage{enumitem}
\usepackage{algorithm}
\usepackage{algpseudocode}
\usepackage{listings}
\usepackage{xcolor}

\hypersetup{
  colorlinks=true,
  linkcolor=blue!50!black,
  citecolor=blue!50!black,
  urlcolor=blue!50!black
}

% ---- Theorem environments ----
\newtheorem{proposition}{Proposition}
\newtheorem{lemma}{Lemma}
\newtheorem{corollary}{Corollary}
\theoremstyle{definition}
\newtheorem{definition}{Definition}
\newtheorem{assumption}{Assumption}
\theoremstyle{remark}
\newtheorem{remark}{Remark}

% ---- Notation shortcuts ----
\newcommand{\R}{\mathbb{R}}
\newcommand{\E}{\mathbb{E}}
\newcommand{\Var}{\operatorname{Var}}
\newcommand{\Cov}{\operatorname{Cov}}
\renewcommand{\Pr}{\mathbb{P}}
\newcommand{\Norm}{\mathcal{N}}
\newcommand{\indicator}{\mathbf{1}}
\newcommand{\given}{\,\vert\,}
\newcommand{\spradius}{\rho_{\!*}}
\DeclareMathOperator*{\argmin}{arg\,min}

% ---- Listings setup for R code ----
\lstset{
  basicstyle=\ttfamily\footnotesize,
  breaklines=true,
  showstringspaces=false,
  language=R,
  commentstyle=\color{gray},
  keywordstyle=\color{blue!60!black},
  stringstyle=\color{green!50!black},
  frame=single,
  framerule=0.3pt
}

% ---- Title block ----
\title{
  A Time-Block Holdout Strategy for Evaluating AR(2)
  Temporal Pooling in a Spatial Skew-$t$ Model\\[4pt]
  \large Research Technical Note: \texttt{time\_block\_strategy}
}
\author{Spatial Skew-$t$ Simulation Study}
\date{\today}

\begin{document}
\maketitle

% ============================================================
\begin{abstract}
A simulation study fits a spatial skew-$t$ model with and without
AR(2) temporal structure on three latent knot-level processes
(precision $\tau$, skewness $z$, and knot location $w$), and scores
predictions with the Brier score. The AR(2) variants fail to dominate
their independent-in-time counterparts, and in one setting score
strictly worse. This note formalizes why: the Brier score is a strictly
proper score for the \emph{marginal} exceedance probability, while the
existing holdout removes spatial \emph{sites} and retains all time
points. Under that design the AR(2) prior and the i.i.d.\ prior induce
the same marginal predictive distribution at an unobserved site, so a
marginal score cannot separate them and the simpler model wins the
bias--variance trade-off. We prove that the conditional $h$-step AR(2)
forecast converges geometrically to the stationary marginal at rate
governed by the companion-matrix spectral radius, define an effective
memory horizon, and use it to design a time-block holdout that
forecasts forward from the posterior of the seam state. We give the
forecasting algorithm, the latent-scale coupling of $\tau$ and $z$, and
a lead-time scoring protocol based on multivariate proper scores. The
result is an evaluation framework under which AR(2) temporal pooling can
genuinely be tested rather than structurally nullified.
\end{abstract}

% ============================================================
\section{Introduction}
\label{sec:intro}

A spatial skew-$t$ model for extreme environmental fields can be
extended with temporal dependence by placing autoregressive priors on
its latent knot-level processes. The natural question is whether this
temporal pooling improves predictive performance for extremes. A
simulation study addressed this by fitting eight method variants ---
including two that activate AR(2) structure on all three latent
processes --- and scoring held-out predictions with the Brier score
across a grid of high threshold quantiles.

The outcome was a null result. The AR(2) variants did not dominate the
independent-in-time skew-$t$ and $t$ variants; in the setting with five
knots the AR(2) model scored \emph{worse} than its own non-temporal
sibling at every threshold, with relative Brier inflation reaching
$1.16$ in the far upper tail. Taken at face value this suggests
temporal pooling is unhelpful. We argue instead that the result is an
artifact of the \emph{interaction between the scoring rule and the
holdout design}, and that the experiment as constructed cannot answer
the question it was built to ask.

This note has three goals. First, in Section~\ref{sec:setup} we
formalize the model, the AR(2) latent processes, and the family of
scoring rules. Second, in Section~\ref{sec:derivations} we prove that a
marginal score under a spatial-only holdout is structurally blind to
temporal dependence, and we quantify the AR(2) memory horizon through
the companion-matrix spectral radius. Third, in
Section~\ref{sec:redesign} we derive a corrected time-block holdout:
its forecasting algorithm, the seam-state conditioning that preserves
the temporal signal, the latent-scale transform order, and a
lead-time scoring protocol. Section~\ref{sec:conclusion} states the
expected outcome and open issues.

% ============================================================
\section{Model Setup and Definitions}
\label{sec:setup}

\subsection{The spatial skew-$t$ observation model}

Let $\mathcal{S} \subset \R^2$ be a study region and let
$t \in \{1,\dots,n_t\}$ index time. Observations are made at sites
$s \in \mathcal{S}$ and times $t$.

\begin{definition}[Observation model]
\label{def:obs}
For site $s$ and time $t$,
\begin{equation}
  y(s,t)
  = x(s,t)^\top \beta
  + \lambda\, z_{\,g(s,t),\,t}
  + \varepsilon(s,t),
  \label{eq:obs}
\end{equation}
where $x(s,t)\in\R^p$ are covariates with coefficient vector
$\beta\in\R^p$; $\lambda$ is a skewness loading;
$\varepsilon(\cdot,t)$ is a mean-zero spatial Gaussian field with
$\Cov\{\varepsilon(s,t),\varepsilon(s',t)\}
= C(s,s')\,\tau_{\,g(s,t),\,t}^{-1}$;
and $C$ is a correlation function with range $\rho$, smoothness $\nu$,
and geometric-anisotropy parameter $\gamma$.
\end{definition}

The model carries $K$ spatial \emph{knots}. For each knot
$k \in \{1,\dots,K\}$ and time $t$ there are three latent quantities:
a precision $\tau_{k,t}>0$, a skewness term $z_{k,t}\ge 0$, and a knot
location $w_{k,t}\in\R^2$. The function $g(s,t)$ returns the index of
the knot nearest to $s$ at time $t$, i.e.\ the Voronoi membership
\begin{equation}
  g(s,t) = \argmin_{k\in\{1,\dots,K\}}
            \big\| s - w_{k,t} \big\|.
  \label{eq:membership}
\end{equation}

\begin{assumption}[Conditional independence of the spatial nugget]
\label{ass:iid-space}
The spatial field $\varepsilon(\cdot,t)$ in \eqref{eq:obs} is drawn
independently across $t$. All cross-time dependence in $y$ is therefore
transmitted exclusively through the latent processes
$\{\tau_{k,t}, z_{k,t}, w_{k,t}\}$.
\end{assumption}

Assumption~\ref{ass:iid-space} is a property of the data-generating
mechanism, not a modelling choice; it localizes the entire temporal
signal in the latent processes and is the hinge of every argument that
follows.

\subsection{Latent AR(2) processes on the Gaussian scale}

Each latent process is generated on a standardized Gaussian scale and
then mapped to its target marginal by a copula transform. Write
$X_t$ generically for any of the latent Gaussian processes
$\tau^\star_{k,\cdot}$, $z^\star_{k,\cdot}$, or $w^\star_{k,\cdot}$.

\begin{definition}[AR(2) latent process]
\label{def:ar2}
The latent Gaussian process obeys
\begin{equation}
  X_t = \phi_1 X_{t-1} + \phi_2 X_{t-2} + \xi_t,
  \qquad \xi_t \overset{\text{iid}}{\sim} \Norm(0,\sigma^2),
  \label{eq:ar2}
\end{equation}
with coefficients $\phi=(\phi_1,\phi_2)$ and innovation variance
$\sigma^2$.
\end{definition}

\begin{definition}[Stationarity triangle]
\label{def:stationary}
The recursion \eqref{eq:ar2} is stationary if and only if
\begin{equation}
  |\phi_2| < 1,
  \qquad \phi_1 + \phi_2 < 1,
  \qquad \phi_2 - \phi_1 < 1.
  \label{eq:triangle}
\end{equation}
\end{definition}

We normalize each latent process to unit marginal variance,
$\gamma_0 \coloneqq \Var(X_t) = 1$. The Yule--Walker relations then fix
the lag-one and lag-two autocovariances and the innovation variance:
\begin{equation}
  \gamma_1 = \frac{\phi_1}{1-\phi_2},
  \qquad
  \gamma_2 = \phi_1\gamma_1 + \phi_2\gamma_0,
  \qquad
  \sigma^2 = \gamma_0 - \phi_1\gamma_1 - \phi_2\gamma_2.
  \label{eq:yw}
\end{equation}

\begin{definition}[Companion form]
\label{def:companion}
Stacking $\bm{X}_t = (X_t, X_{t-1})^\top$, the recursion
\eqref{eq:ar2} becomes the first-order vector system
\begin{equation}
  \bm{X}_t = F\,\bm{X}_{t-1} + \bm{e}_t,
  \qquad
  F = \begin{pmatrix} \phi_1 & \phi_2 \\ 1 & 0 \end{pmatrix},
  \qquad
  \bm{e}_t = \begin{pmatrix} \xi_t \\ 0 \end{pmatrix},
  \label{eq:companion}
\end{equation}
with innovation covariance $Q = \operatorname{diag}(\sigma^2, 0)$.
Stationarity \eqref{eq:triangle} is equivalent to the spectral radius
condition $\spradius \coloneqq \rho(F) < 1$, where $\rho(F)$ is the
modulus of the largest eigenvalue of $F$.
\end{definition}

\begin{definition}[Stationary state distribution]
\label{def:statedist}
The stationary distribution of the companion state is
$\bm{X}_t \sim \Norm(\bm{0}, \Sigma)$, where $\Sigma$ solves the
discrete Lyapunov equation
\begin{equation}
  \Sigma = F\,\Sigma\,F^\top + Q,
  \qquad
  \Sigma =
  \begin{pmatrix} \gamma_0 & \gamma_1 \\ \gamma_1 & \gamma_0 \end{pmatrix}
  =
  \begin{pmatrix} 1 & \gamma_1 \\ \gamma_1 & 1 \end{pmatrix}.
  \label{eq:lyapunov}
\end{equation}
\end{definition}

Initializing a simulated path by drawing $(X_1,X_2)$ from
$\Norm(\bm 0,\Sigma)$ yields an \emph{exactly} stationary series with no
burn-in. This is the construction used by the study's generator: a
manual AR($p$) recursion with an exact stationary initial state.

\subsection{Copula transforms to observable scale}

The latent Gaussian processes are mapped to their target marginals by
monotone copula transforms, applied \emph{after} the AR(2) recursion:
\begin{align}
  \tau_{k,t}   &= G^{-1}_{\alpha,\beta}\!\big(\Phi(\tau^\star_{k,t})\big),
  \label{eq:tau-transform}\\
  z_{k,t}      &= H^{-1}_{\sigma_z}\!\big(\Phi(z^\star_{k,t})\big),
  \qquad \sigma_z = \tau_{k,t}^{-1/2},
  \label{eq:z-transform}\\
  w_{k,t}      &= \Psi^{-1}\!\big(\Phi(w^\star_{k,t})\big),
  \label{eq:w-transform}
\end{align}
where $\Phi$ is the standard normal CDF, $G^{-1}_{\alpha,\beta}$ is a
gamma quantile function, $H^{-1}_{\sigma_z}$ is a half-normal quantile
function with scale $\sigma_z$, and $\Psi^{-1}$ is a scaled inverse-probit
map onto the spatial domain. Equation~\eqref{eq:z-transform} records a
structural coupling: the skewness scale $\sigma_z$ is itself a function
of the precision $\tau_{k,t}$.

\subsection{Scoring rules}
\label{sec:scores}

Let $F$ denote a predictive distribution and $y$ the realized outcome.

\begin{definition}[Brier score]
\label{def:brier}
For threshold $c$, with predicted exceedance probability
$\hat p(s,t;c) = \Pr_F\{Y(s,t) > c\}$ and indicator
$\indicator\{y(s,t)>c\}$, the Brier score is the cell-averaged
quadratic loss
\begin{equation}
  \mathrm{BS}(c)
  = \frac{1}{|\mathcal{V}|}
    \sum_{(s,t)\in\mathcal{V}}
    \big(\indicator\{y(s,t)>c\} - \hat p(s,t;c)\big)^2,
  \label{eq:brier}
\end{equation}
where $\mathcal{V}$ is the set of held-out $(s,t)$ cells.
\end{definition}

\begin{definition}[Energy score]
\label{def:energy}
For a vector outcome $\bm y \in \R^d$ and predictive $F$ on $\R^d$,
\begin{equation}
  \mathrm{ES}(F,\bm y)
  = \E_F \|\bm X - \bm y\|
  - \tfrac12\, \E_F \|\bm X - \bm X'\|,
  \label{eq:energy}
\end{equation}
where $\bm X,\bm X' \overset{\text{iid}}{\sim} F$. The energy score is a
strictly proper score for the \emph{joint} law of $\bm y$. Its
univariate case $d=1$ is the continuous ranked probability score
(CRPS).
\end{definition}

\begin{definition}[Variogram score of order $p$]
\label{def:variogram}
For $\bm y \in \R^d$ and weights $w_{ij}\ge 0$,
\begin{equation}
  \mathrm{VS}_p(F,\bm y)
  = \sum_{i=1}^d \sum_{j=1}^d
    w_{ij}
    \Big(
      |y_i - y_j|^p - \E_F |X_i - X_j|^p
    \Big)^2 .
  \label{eq:variogram}
\end{equation}
Indexing $i,j$ by time lag makes $\mathrm{VS}_p$ sensitive to the
autocorrelation structure of the predictive.
\end{definition}

% ============================================================
\section{Why a Marginal Score Cannot See Temporal Structure}
\label{sec:derivations}

This section establishes the diagnosis. The argument has three steps:
the Brier score depends on the predictive only through marginals
(Proposition~\ref{prop:brier-marginal}); a spatial-only holdout makes
the AR(2) and i.i.d.\ models share those marginals
(Proposition~\ref{prop:collapse}); and the AR(2) advantage that
\emph{does} exist lives in a forecast horizon that the holdout never
exercises (Propositions~\ref{prop:forecast}--\ref{prop:horizon}).

\subsection{The Brier score is a marginal functional}

\begin{proposition}[Brier insensitivity to dependence]
\label{prop:brier-marginal}
The Brier score \eqref{eq:brier} depends on the predictive
distribution $F$ only through the collection of univariate marginal
exceedance probabilities $\{\hat p(s,t;c) : (s,t)\in\mathcal{V}\}$.
Consequently, if two predictive models $F_A$ and $F_B$ induce identical
marginals at every held-out cell --- $\hat p_A(s,t;c) = \hat p_B(s,t;c)$
for all $(s,t)\in\mathcal{V}$ and all $c$ --- then they receive
identical Brier scores, regardless of any difference in their
dependence structure.
\end{proposition}

\begin{proof}
Each summand of \eqref{eq:brier} is a function of the pair
$\big(\indicator\{y(s,t)>c\},\,\hat p(s,t;c)\big)$. The indicator is a
property of the data, not of $F$. The only term through which $F$
enters is the scalar $\hat p(s,t;c) = \Pr_F\{Y(s,t)>c\}$, which is a
functional of the univariate marginal of $F$ at cell $(s,t)$. The sum
contains no term coupling two distinct cells. Hence $\mathrm{BS}(c)$ is
a function of the marginals alone, and any two predictives agreeing on
all marginals produce the same value. \qed
\end{proof}

\begin{remark}
\label{rem:proper}
Proposition~\ref{prop:brier-marginal} is not a defect of the Brier
score; it is its design. The Brier score is a strictly proper score for
a \emph{marginal} binary event. Cross-cell dependence is simply outside
the quantity it was built to assess. The energy and variogram scores of
Definitions~\ref{def:energy}--\ref{def:variogram} are the proper scores
that \emph{do} interrogate dependence.
\end{remark}

\subsection{Spatial-only holdout collapses the comparison}

Consider the holdout used by the study: a subset of spatial sites
$\mathcal{S}_{\text{val}}$ is withheld, while \emph{every} time point is
observed at the training sites. Prediction targets are the cells
$\{(s,t): s\in\mathcal{S}_{\text{val}},\ t = 1,\dots,n_t\}$.

\begin{proposition}[Structural collapse under spatial-only holdout]
\label{prop:collapse}
Under Assumption~\ref{ass:iid-space} and a spatial-only holdout, the
predictive marginal at an unobserved site $s^\star$ and time $t$ is
\begin{equation}
  \Pr\{Y(s^\star,t) > c\}
  = \E_{\,\tau,z,w}\!\left[
      \Pr\big\{ Y(s^\star,t) > c \,\big|\, \tau_{\cdot,t}, z_{\cdot,t},
                w_{\cdot,t} \big\}
    \right],
  \label{eq:marginal-pred}
\end{equation}
where the expectation is over the time-$t$ \emph{marginal} law of the
latent processes. The AR(2) prior and the i.i.d.\ prior, when both are
normalized to the same stationary marginal $\Norm(0,1)$ on the latent
scale, induce the same law for $(\tau_{\cdot,t}, z_{\cdot,t},
w_{\cdot,t})$ at each fixed $t$, hence the same marginal
\eqref{eq:marginal-pred}, hence --- by
Proposition~\ref{prop:brier-marginal} --- the same Brier score in
expectation.
\end{proposition}

\begin{proof}
By Assumption~\ref{ass:iid-space} the spatial nugget at time $t$ is
independent of all other times, so conditioning on the time-$t$ latent
slice $(\tau_{\cdot,t}, z_{\cdot,t}, w_{\cdot,t})$ renders
$Y(s^\star,t)$ independent of every other time. The exceedance
probability at cell $(s^\star,t)$ therefore depends on the latent
processes only through their time-$t$ values, giving
\eqref{eq:marginal-pred}.

The AR(2) prior of Definition~\ref{def:ar2} and an i.i.d.\ prior differ
only in the \emph{joint} law across $t$; both are constructed with the
same unit-variance stationary marginal $X_t \sim \Norm(0,1)$ on the
latent scale (the normalization \eqref{eq:yw}). The copula transforms
\eqref{eq:tau-transform}--\eqref{eq:w-transform} are applied
marginally and time-pointwise, so equal latent marginals yield equal
marginals for $(\tau_{\cdot,t},z_{\cdot,t},w_{\cdot,t})$. The
right-hand side of \eqref{eq:marginal-pred} is thus identical for the
two priors. Proposition~\ref{prop:brier-marginal} converts the equal
marginals into equal Brier scores. \qed
\end{proof}

\begin{remark}
\label{rem:bias-variance}
Proposition~\ref{prop:collapse} explains the empirical observation. If
the AR(2) and i.i.d.\ models are scored on a quantity for which they
are exchangeable, the comparison reduces to a bias--variance contest.
The AR(2) model carries six extra parameters --- a pair
$(\phi_1,\phi_2)$ for each of $\tau$, $z$, $w$ --- whose posterior
uncertainty inflates the predictive variance without changing the
scored marginal. The simpler model therefore wins, and the AR(2)
penalty grows in the upper tail where estimation variance is largest.
This is a property of the experiment, not of temporal pooling.
\end{remark}

\subsection{Where the AR(2) advantage actually lives}

The AR(2) advantage is a statement about \emph{conditional} forecasts.
We make it precise.

\begin{proposition}[Conditional $h$-step forecast]
\label{prop:forecast}
Let $T_o$ be the last observed time and $\bm X_{T_o}$ the companion
state \eqref{eq:companion}. For a stationary AR(2),
\begin{equation}
  \E[\bm X_{T_o+h} \given \bm X_{T_o}] = F^h\,\bm X_{T_o},
  \qquad
  \Var[\bm X_{T_o+h} \given \bm X_{T_o}]
    = \Sigma_h \coloneqq \sum_{j=0}^{h-1} F^j Q\,(F^j)^\top .
  \label{eq:hstep}
\end{equation}
As $h\to\infty$, $F^h \to \bm 0$ and $\Sigma_h \to \Sigma$, the
stationary covariance of Definition~\ref{def:statedist}.
\end{proposition}

\begin{proof}
Iterating \eqref{eq:companion} forward $h$ steps gives
$\bm X_{T_o+h} = F^h \bm X_{T_o}
  + \sum_{j=0}^{h-1} F^j \bm e_{T_o+h-j}$.
The innovations $\bm e_{\,\cdot}$ are mean-zero and independent of
$\bm X_{T_o}$, which yields the conditional mean and variance in
\eqref{eq:hstep}. Stationarity gives $\spradius = \rho(F) < 1$
(Definition~\ref{def:companion}), so $F^h \to \bm 0$ geometrically;
the partial sum $\Sigma_h$ is monotone increasing in the positive
semidefinite order and converges to the solution $\Sigma$ of the
Lyapunov equation \eqref{eq:lyapunov}. \qed
\end{proof}

\begin{corollary}[The stationary distribution is the infinite-horizon limit]
\label{cor:stationary-limit}
Drawing the forecast origin from the stationary distribution
$\Norm(\bm 0,\Sigma)$ is equivalent to the limit $h\to\infty$ of
\eqref{eq:hstep}. It discards the conditioning information in
$\bm X_{T_o}$ and reproduces, for every $h$, the i.i.d.\ model's
marginal predictive.
\end{corollary}

\begin{proof}
Setting the mean term $F^h\bm X_{T_o}$ to zero and the variance to
$\Sigma$ in \eqref{eq:hstep} is exactly the result of replacing
$\bm X_{T_o}$ by a draw from $\Norm(\bm 0,\Sigma)$ and recursing: the
forecast law becomes $\Norm(\bm 0,\Sigma)$ for all $h$, independent of
$h$. That is the time-$t$ stationary marginal, which by
Proposition~\ref{prop:collapse} is the i.i.d.\ predictive. \qed
\end{proof}

Corollary~\ref{cor:stationary-limit} is the central implementation
warning: initializing a forecast from the stationary distribution
\emph{is} the act of erasing the AR(2) signal. The forecast must
instead be conditioned on the posterior of the seam state $\bm X_{T_o}$.

\begin{proposition}[Effective memory horizon]
\label{prop:horizon}
The discrepancy between the conditional forecast \eqref{eq:hstep} and
the stationary marginal decays geometrically:
\begin{equation}
  \big\| \E[\bm X_{T_o+h}\given \bm X_{T_o}] \big\|
    \le \|F^h\|\,\|\bm X_{T_o}\|
    = O\!\big(\spradius^{\,h}\big),
  \qquad
  \|\Sigma - \Sigma_h\| = O\!\big(\spradius^{\,2h}\big).
  \label{eq:decay}
\end{equation}
Hence for a tolerance $\epsilon>0$ there is an effective memory horizon
\begin{equation}
  h^\star(\epsilon)
  = \left\lceil \frac{\log\epsilon}{\log\spradius} \right\rceil,
  \label{eq:hstar}
\end{equation}
beyond which the AR(2) and i.i.d.\ predictives are indistinguishable to
within $\epsilon$.
\end{proposition}

\begin{proof}
Since $\spradius<1$, the spectral radius formula gives
$\|F^h\| = O(\spradius^{\,h})$ up to a polynomial factor absorbed into
the constant; the mean bound in \eqref{eq:decay} follows. For the
variance, $\Sigma - \Sigma_h = \sum_{j\ge h} F^j Q (F^j)^\top$, whose
norm is bounded by a geometric tail of ratio $\spradius^{\,2}$, giving
$O(\spradius^{\,2h})$. Setting the mean bound equal to $\epsilon$ and
solving for $h$ yields \eqref{eq:hstar}. \qed
\end{proof}

\begin{remark}
\label{rem:horizon-numbers}
Proposition~\ref{prop:horizon} converts a modelling intuition into a
design parameter. For coefficients of the order
$\phi=(0.6,-0.3)$ the spectral radius is moderate and $h^\star$ is a
single-digit number of steps. A holdout that scores predictions only
beyond $h^\star$ measures nothing but the stationary marginal; a
holdout that never forecasts past $T_o$ at all --- the spatial-only
design --- exercises the conditional forecast for exactly zero steps.
Either way the AR(2) signal is invisible by construction.
\end{remark}

% ============================================================
\section{A Time-Block Holdout That Exercises Temporal Pooling}
\label{sec:redesign}

The diagnosis dictates the cure: hold out \emph{time blocks} and
forecast forward from the posterior of the seam state, scoring with a
rule that credits joint structure and resolving the result by lead
time.

\subsection{Holdout geometry}

\begin{definition}[Time-block holdout]
\label{def:blocks}
Choose $B$ non-overlapping blocks indexed $b=1,\dots,B$, each of length
$H$ with $H \lesssim h^\star$ from \eqref{eq:hstar} (a practical range
is $H\in[5,15]$). Block $b$ has seam time $T_o^{(b)}$ and prediction
targets $\{(s,t): s\in\mathcal{S},\ t = T_o^{(b)}+1,\dots,T_o^{(b)}+H\}$.
Blocks are spread across the record so adjacent blocks are
approximately independent.
\end{definition}

Choosing $H$ near $h^\star$ keeps every scored lead time inside the
window where the conditional forecast still differs from the stationary
marginal (Proposition~\ref{prop:horizon}).

\subsection{Forecasting algorithm}

The forecast must satisfy three requirements derived above:
(i) condition on the posterior of the seam state, not the stationary
distribution (Corollary~\ref{cor:stationary-limit});
(ii) propagate full parameter \emph{and} state uncertainty by iterating
over posterior draws; and
(iii) recurse on the latent Gaussian scale, applying the copula
transforms \eqref{eq:tau-transform}--\eqref{eq:w-transform} only after
the recursion.

\begin{algorithm}[h]
\caption{Posterior predictive forecast over a held-out time block}
\label{alg:forecast}
\begin{algorithmic}[1]
\State \textbf{Input:} MCMC draws $m=1,\dots,M$; seam time $T_o$;
        horizon $H$; held-out sites $\mathcal{S}_{\text{val}}$.
\State Set a single RNG seed \emph{once}, outside all loops.
\For{each draw $m = 1,\dots,M$}
  \For{each latent process $X \in \{\tau^\star, z^\star, w^\star\}$
        and knot $k$}
    \State Extract seam state
           $\big(X^{(m)}_{k,T_o-1},\, X^{(m)}_{k,T_o}\big)$
           directly from draw $m$'s imputed latent trajectory.
    \For{$h = 1,\dots,H$}
      \State $X^{(m)}_{k,T_o+h} \gets
              \phi_1^{(m)} X^{(m)}_{k,T_o+h-1}
            + \phi_2^{(m)} X^{(m)}_{k,T_o+h-2}
            + \sigma^{(m)}\,\xi$,
            \quad $\xi\sim\Norm(0,1)$.
    \EndFor
  \EndFor
  \For{$h = 1,\dots,H$}
    \State $\tau^{(m)}_{\cdot,T_o+h} \gets
            G^{-1}_{\alpha^{(m)},\beta^{(m)}}
            \!\big(\Phi(\tau^{\star(m)}_{\cdot,T_o+h})\big)$.
    \State $z^{(m)}_{\cdot,T_o+h} \gets
            H^{-1}_{\sigma_z}\!\big(\Phi(z^{\star(m)}_{\cdot,T_o+h})\big)$,
            \quad $\sigma_z = (\tau^{(m)}_{\cdot,T_o+h})^{-1/2}$.
    \State $w^{(m)}_{\cdot,T_o+h} \gets
            \Psi^{-1}\!\big(\Phi(w^{\star(m)}_{\cdot,T_o+h})\big)$;
            recompute membership $g$ by \eqref{eq:membership}.
    \For{each $s \in \mathcal{S}_{\text{val}}$}
      \State $\mu \gets x(s,T_o+h)^\top\beta^{(m)}
              + \lambda^{(m)} z^{(m)}_{g(s,T_o+h),T_o+h}$.
      \State Draw $y^{\star(m)}(s,T_o+h)$ from the spatial Gaussian
              field with mean $\mu$ and precision
              $\tau^{(m)}_{g(s,T_o+h),T_o+h}$.
    \EndFor
  \EndFor
\EndFor
\State \textbf{Output:} predictive sample
        $\{y^{\star(m)}(s,T_o+h)\}_{m,s,h}$.
\end{algorithmic}
\end{algorithm}

\begin{proposition}[Coherence of Algorithm~\ref{alg:forecast}]
\label{prop:coherence}
Algorithm~\ref{alg:forecast} produces draws from the posterior
predictive distribution of the time block, with both parameter
uncertainty and seam-state uncertainty preserved, and with the
$\tau$--$z$ coupling \eqref{eq:z-transform} respected.
\end{proposition}

\begin{proof}
For each $m$, the tuple
$\big(\phi^{(m)}, \sigma^{(m)}, \bm X^{(m)}_{T_o},
       \beta^{(m)}, \lambda^{(m)}, \alpha^{(m)}, \beta^{(m)},
       \rho^{(m)}, \nu^{(m)}, \gamma^{(m)}\big)$
is a single coherent sample from the joint posterior: the latent state
$\bm X^{(m)}_{T_o}$ is an imputed quantity of draw $m$, not a fresh
draw, so by Proposition~\ref{prop:forecast} the recursion forecasts
from the correct conditional law and, by
Corollary~\ref{cor:stationary-limit}, does not collapse to the
stationary marginal. Averaging the predictive over $m=1,\dots,M$
integrates the posterior; no plug-in of posterior means occurs, so
parameter uncertainty is retained. The $\tau$--$z$ coupling is enforced
in line 12: the half-normal scale $\sigma_z$ is evaluated at draw $m$'s
own forecasted $\tau^{(m)}_{\cdot,T_o+h}$, so $\tau$ and $z$ are
transformed jointly within the draw rather than from independent
marginals. The AR innovations of $\tau^\star$ and $z^\star$ are a priori
independent (Definition~\ref{def:ar2}), so no further cross-process
innovation term is required; the posterior coupling is transmitted
entirely through the shared seam state and shared parameters. \qed
\end{proof}

\begin{remark}[Implementation pitfalls]
\label{rem:pitfalls}
Four failure modes void Proposition~\ref{prop:coherence} in practice.
(a) \emph{Seam from the stationary distribution} --- redrawing
$\bm X_{T_o}$ from $\Norm(\bm 0,\Sigma)$ instead of extracting it from
draw $m$; by Corollary~\ref{cor:stationary-limit} this erases the
signal.
(b) \emph{Plug-in means} --- forecasting from posterior means of
$\phi$ and $\bm X_{T_o}$; this collapses uncertainty and under-disperses
the predictive.
(c) \emph{Reseeding inside the loop} --- any \texttt{set.seed} call
within the $m$- or $h$-loop breaks path coherence; the seed is set once,
outside.
(d) \emph{Non-stationary draws} --- posterior tail draws of $\phi$ may
exit the triangle \eqref{eq:triangle}, making the recursion explode;
either reparameterize to partial autocorrelations or reject such draws
before forecasting.
A further numerical hazard is Voronoi membership flipping: small drift
in the forecasted $w_{k,t}$ can move a site across a knot boundary and
discontinuously change $\mu$. Averaging the predictive over many draws
softens the partition; reporting a fixed-membership ablation
($g(s,t)\equiv g(s,T_o)$) alongside the full result isolates the AR(2)
signal in $\tau$ and $z$ from partition noise.
\end{remark}

\subsection{Lead-time scoring protocol}

The energy and variogram scores of
Definitions~\ref{def:energy}--\ref{def:variogram} are multivariate:
each returns one number for an entire block vector and cannot be split
by lead time. A lead-time curve therefore requires a \emph{univariate}
proper score evaluated separately at each $h$.

\begin{definition}[Lead-time score curve]
\label{def:leadcurve}
Let $S\big(h; b, s\big)$ be a univariate proper score (e.g.\ CRPS, or
the Brier score restricted to lead $h$) for the prediction at site $s$,
lead time $h$, in block $b$. The lead-time curve is
\begin{equation}
  \bar S(h)
  = \frac{1}{B\,|\mathcal{S}_{\text{val}}|}
    \sum_{b=1}^{B}
    \sum_{s\in\mathcal{S}_{\text{val}}}
    S\big(h; b, s\big),
  \qquad h = 1,\dots,H,
  \label{eq:leadcurve}
\end{equation}
with a standard error at each $h$ estimated from the dispersion of the
$B$ per-block means.
\end{definition}

The evaluation reports two complementary objects.
\begin{enumerate}[label=(\roman*)]
  \item \textbf{Lead-time curve.} Plot $\bar S(h)$ from
        \eqref{eq:leadcurve} for the AR(2) model and the i.i.d.\
        baseline on common axes, with error bands. By
        Proposition~\ref{prop:horizon} the AR(2) curve should lie below
        the baseline at small $h$ and converge to it near
        $h^\star$. The crossing point is the headline quantity.
  \item \textbf{Joint-structure summary.} Report the energy score
        \eqref{eq:energy} and variogram score \eqref{eq:variogram}
        computed over the per-site length-$H$ time vector, averaged
        over blocks and sites --- one number per model that credits the
        temporal dependence the lead-time curve resolves.
\end{enumerate}

\begin{remark}[Verification checks]
\label{rem:checks}
Three diagnostics confirm the forecast loop is correct.
(1) With $H=0$, re-scoring time $T_o$ must give AR(2) and i.i.d.\ the
same predictive --- any difference is a seam-indexing bug.
(2) A draw with $\phi=\bm 0$ must collapse immediately to the marginal;
residual AR-like behaviour indicates double-counted state.
(3) A long horizon $H\gg h^\star$ with stationary $\phi$ must reproduce
the $\Norm(0,1)$ latent marginal across draws, validating the
innovation-variance normalization \eqref{eq:yw}.
\end{remark}

% ============================================================
\section{Conclusion and Future Work}
\label{sec:conclusion}

The null result --- AR(2) temporal pooling failing to beat an
i.i.d.\ baseline on the Brier score --- is fully explained by
Propositions~\ref{prop:brier-marginal} and \ref{prop:collapse}: a
marginal score under a spatial-only holdout makes the two models
exchangeable, so the comparison degenerates to a bias--variance contest
that the lower-dimensional model wins. The finding is a property of the
experimental design, not evidence against temporal pooling.

The corrected framework of Section~\ref{sec:redesign} removes the
structural blindness. By holding out time blocks of length $H\lesssim
h^\star$, forecasting from the posterior of the seam state
(Algorithm~\ref{alg:forecast}, Proposition~\ref{prop:coherence}), and
resolving performance by lead time \eqref{eq:leadcurve}, the evaluation
exercises precisely the conditional forecast in which the AR(2)
advantage resides (Proposition~\ref{prop:forecast}). The expected
signature of a genuine effect is a lead-time curve on which AR(2) beats
the i.i.d.\ baseline at short horizons and converges to it near
$h^\star$ (Proposition~\ref{prop:horizon}).

Open items for the next iteration:
\begin{itemize}[leftmargin=1.4em]
  \item \textbf{Block-length calibration.} Estimate $\spradius$ from
        the posterior of $\phi$ and set $H$ from \eqref{eq:hstar} per
        setting rather than by a fixed rule.
  \item \textbf{Partition uncertainty.} Compare the soft-membership
        predictive against the fixed-membership ablation of
        Remark~\ref{rem:pitfalls}; report both so the AR(2) signal in
        $\tau,z$ is separated from Voronoi flipping noise.
  \item \textbf{Score sensitivity.} Contrast the lead-time CRPS curve
        with the energy and variogram summaries to confirm the
        conclusions are not an artifact of a single scoring rule.
  \item \textbf{Forecasting deployment.} Beyond block holdout, evaluate
        genuine forward forecasting from the end of the record, the
        setting in which temporal pooling is of direct operational
        value.
\end{itemize}

% ============================================================
\begin{thebibliography}{9}

\bibitem{gneiting2007}
T.~Gneiting and A.~E.~Raftery,
\emph{Strictly Proper Scoring Rules, Prediction, and Estimation},
Journal of the American Statistical Association, 102(477):359--378, 2007.

\bibitem{scoringrules}
A.~Jordan, F.~Kr\"uger, and S.~Lerch,
\emph{Evaluating Probabilistic Forecasts with \texttt{scoringRules}},
Journal of Statistical Software, 90(12):1--37, 2019.

\bibitem{brockwell}
P.~J.~Brockwell and R.~A.~Davis,
\emph{Time Series: Theory and Methods}, 2nd ed.,
Springer, New York, 1991.

\bibitem{scheuerer2015}
M.~Scheuerer and T.~M.~Hamill,
\emph{Variogram-Based Proper Scoring Rules for Probabilistic Forecasts
of Multivariate Quantities},
Monthly Weather Review, 143(4):1321--1334, 2015.

\end{thebibliography}

\end{document}
